% ============================================
% LUC-19 — A Última Noite em Casa (edição revista R1075)
% ============================================
\fragtitle{luc}{L}{LUC}{19}{A Última Noite em Casa}
\markboth{LUC-19 --- A Última Noite em Casa}{}

\noindent \textit{Belo Horizonte --- 13 de maio de 2026, 23h48}

\noindent Lúcia fez as malas como quem organiza um arquivo. Três blusas, duas calças, roupa de baixo, o carregador, o notebook, o estojo com os documentos. Não levou livros. Não levou fotografias. Não levou o perfume. Havia uma lógica nisso, embora ela não conseguisse explicar: tudo o que ficasse no apartamento seria encontrado por alguém. Tudo o que fosse com ela seria perdido com ela.

\noindent Na estante, entre os manuais de psiquiatria forense, havia uma caixa de sapatos com fotos antigas. Lúcia abriu. A primeira era dela aos oito anos, na praia de Santos, com o pai. Ele estava de short, camisa aberta, sorrindo para a câmera. Ela tinha os olhos apertados por causa do sol. Havia areia no pé dela. O pai morreria doze anos depois, mas naquela foto ainda tinha todo o tempo do mundo.

\noindent Lúcia olhou para a foto por um minuto. Depois colocou na caixa, fechou, guardou na estante. Não levou.

\noindent O celular vibrou. Mensagem da mãe: ``Você vem domingo?''. Lúcia digitou ``Vou''. Apagou. Digitou ``Acho que não''. Apagou. Digitou ``Te amo''. Enviou. Depois desligou o telefone.

\noindent Havia uma garrafa de vinho aberta na cozinha, pela metade. Lúcia serviu um copo, bebeu em pé, olhando para a parede. O vinho tinha gosto de ferro. Ela colocou o copo na pia, lavou, secou, guardou no armário. Não queria deixar nada fora do lugar.

\noindent Sentou-se à mesa com uma folha de papel e uma caneta. Escreveu para Ronaldo:

\noindent ``Ronaldo, não fui eu que escolhi isso. Você também não. Mas nós dois empurramos. Se eu não voltar, não procure. Não abra a caixa. Não leia o que eu escrevi. E não conte para a minha mãe o que encontrou. Deixe que ela ache que foi um acidente. É mais gentil.''

\noindent Leu a carta. A letra era a dela. Desta vez, o ``R'' estava normal. Rasgou a folha ao meio. Jogou no lixo. Não queimou. Não havia mais fósforos.

\noindent Às 23h58, o interfone tocou. Lúcia não esperava ninguém. Atendeu.

\noindent ``Entrega para Lúcia Mendes.''

\noindent ``Quem?''

\noindent ``Não tem nome. É um envelope.''

\noindent Ela desceu. Na portaria, havia um envelope pardo sobre o balcão. O porteiro não se lembrava de quem entregou. Lúcia abriu. Dentro, uma única folha de papel, sem cabeçalho, com uma frase datilografada:

\noindent ``A página 62 está esperando.''

\noindent Ela subiu. No apartamento, a caixa de metal estava aberta. O diário, dentro, aberto na página 62. Lúcia não o tinha deixado aberto. Tinha certeza. Fechou os olhos. Contou até sessenta e dois. Abriu.

\noindent A página estava em branco. Não completamente. No canto inferior direito, havia uma linha escrita a lápis, com letra infantil: ``Eu não quero ir sozinho.''

\noindent Lúcia sentou no chão. Ficou olhando para a frase. Havia uma tristeza naquilo que não era horror. Era a tristeza de uma criança que não quer entrar no escuro. Ela pensou em Joaquim, aos oito anos, na estrada de Senador Pompeu, segurando a barra da saia da mãe. Pensou no pai dele, sentado debaixo do juazeiro. Pensou em si mesma, aos oito anos, na praia de Santos, com o pai ainda vivo. Havia uma linha entre aquelas duas crianças, e a linha era o Brasil.

\noindent ``Eu vou com você'', ela disse.

\noindent Não sabia a quem estava respondendo. Não importava.

\noindent Guardou o diário na caixa. Fechou o cadeado. Colocou a chave no copo com água. Vestiu o casaco. Pegou a mala. Antes de sair, olhou para o apartamento uma última vez: a sala pequena, a estante, a mesa, a foto na parede. Não sentiu despedida. Sentiu que já tinha saído havia muito tempo.

\noindent Na portaria, o porteiro dormia. Lúcia colocou a chave do apartamento sobre o balcão. Saiu. A rua estava vazia. O carro estava na vaga, com o para-brisa molhado. Ela abriu a porta, colocou a mala no banco de trás, sentou-se. Antes de ligar o motor, olhou para o copo que ficara na mesa, lá em cima, com a chave dentro.

\noindent A água estava morna. Ela sabia.

> **↪ Rotas:** **LUC-18** · **LUC-20** · **CONT-01**
